% problem-000343 7 碗形芳香二价阴离子
\section*{7. 碗形芳香二价阴离子}
B在质谱仪中的⾃由区场中寿命约为1微秒，在常温下不能分离得到。三⼗年后化学家们终于由A合成了第⼀个碗形芳⾹⼆价阴离⼦ ${ \bf C } , ~ [ { \bf C } _ { 1 0 } \mathrm { H } _ { 6 } ] ^ { 2 - }$ 。化合物C中六个氢的化学环境相同，在⼀定条件下可以转化为 $\mathbf { B } _ { \circ }$ 。化合物A转化为C的过程如下所⽰：

$
\begin{array}{c} \mathrm{C} _ {1 0} \mathrm{H} _ {1 0} \xrightarrow [ \text {hexane} ]{\mathrm{nBuLi,} ^ {\mathrm{t}} \mathrm{BuOK,} ^ {\mathrm{n}} \mathrm{C} _ {6} \mathrm{H} _ {1 4}} [ \mathrm{C} _ {1 0} \mathrm{H} _ {6} ] ^ {2 -} \cdot 2 \mathrm{K} ^ {+} \xrightarrow [ \text {Et} _ {2} \mathrm{O} , \text {hexane} ]{\mathrm{Me} _ {3} \mathrm{SnX}} \xrightarrow [ \text {-78°C} ]{\mathrm{MeLi}} [ \mathrm{C} _ {1 0} \mathrm{H} _ {6} ] ^ {2 -} \cdot 2 \mathrm{Li} ^ {+} \\ \mathbf {A} \end{array}
$

\subsection*{7-1}
A 的结构简式：

\subsection*{7-2}
B的结构简式：

\subsection*{7-3}
C的结构简式：

\subsection*{7-4}
B是否具有芳⾹性？为什么？
